\par\SourceLeaf{165}
Sollen die Änderungen für jetzt aber auf das äußerste zulässige Minimum beschränkt bleiben, so ist die Kommission der Meinung, dass dies auf möglichst schonende Weise durch den folgenden Plan geschehen kann, den sie deshalb unter dieser Bedingung zur Annahme empfiehlt.

I.\quad Die jährlichen Beiträge werden auf 3 Louis d'or erhöht.

II.\quad Die Pensionen bleiben für jetzt bei 250 Rthl., solange die Zahl der Witwen\footnote{Der Kürze wegen werden hier immer nur Witwen genannt; es versteht sich aber von selbst, dass stets eine Waisenpension wie eine Witwenpension gezählt werden muss.} ohne die Professorenwitwe G. nicht über 18 hinausgeht; dies ist die gegenwärtige Anzahl. Andernfalls wird die Pension so festgesetzt, dass zum festen Teil von 200 Rthl. noch die Dividende hinzukommt, die auf jede einzelne Witwe entfällt, wenn man 900 Rthl. (als 18 mal 50 Rthl.) gleichmäßig unter die vorhandenen Witwen verteilt. Bei der höchsten bisher vorgekommenen Witwenzahl, nämlich 21 (ohne die Professorenwitwe G.), würde die Pension also noch 242 Rthl. 20 Ggr. 7 Pf. Gold betragen.

Man erkennt leicht, dass das Defizit bei dieser Einrichtung noch nicht vollständig gedeckt ist: Tatsächlich ergibt die Rechnung, dass, wenn der Beginn der Wirksamkeit vom 1. Oktober 1845 an vorausgesetzt wird, noch 2963 Rthl. ungedeckt bleiben. Beginnt der erhöhte Beitrag später, bleibt also verhältnismäßig mehr ungedeckt. Wenn jedoch der gegenwärtig noch bestehende höhere Zinsertrag vorerst ohne erhebliche Schmälerung anhält und sonst keine bedeutenden Verluste eintreten, könnte man hoffen, dass jener Defizitrest nach einer mäßigen Anzahl von Jahren von selbst gedeckt sein wird. Nimmt die Gesellschaft diesen Plan an, darf sie sich jedoch nicht verhehlen, dass die Witwenkasse bei einer so knappen Bemessung der Mittel im Verhältnis zu den Verpflichtungen einem schwer beladenen Schiff gleicht, das in seichtem Fahrwasser fährt und seine Sicherheit nur im fortwährenden wachsamen Loten findet.

Die Kommission hält es daher für unumgänglich notwendig, dass mit der obigen Regelung noch verbunden wird

III.\quad eine sorgfältige neue Bilanzrechnung, die in angemessenen Zeitabständen (etwa alle 10 oder 5 Jahre) zu wiederholen ist, und dass je nach ihren Ergebnissen gegebenenfalls weitere Abhilfe vorbehalten bleibt. Als notwendige vorbereitende Maßnahme wäre von jetzt an eine genaue Buchführung über die für die Witwenkasse maßgeblichen persönlichen Verhältnisse aller Beteiligten einzuführen, wie bereits an mehreren Stellen der Denkschrift angedeutet.

Die Kommission hält es nicht für nötig, die Vorteile dieser Regelung hier ausführlich darzulegen, und macht nur auf Folgendes aufmerksam. Zeigt sich nach der ersten oder zweiten Revision, dass sich die Bilanz nicht nur verbessert, sondern viel stärker verbessert hat, als man hoffen konnte, so wird man die Pensionen weiter von 250 Rthl. auf 260 Rthl. erhöhen dürfen, falls die Witwenzahl unter 18 liegt, oder bei hoher Witwenzahl die gesamte Zusatzdividende von 900 Rthl. auf 1080 Rthl. (das heißt 18 mal 60 Rthl.). Zeigt sich bei der Revision dagegen eine Verschlechterung der Bilanz, wird man sich einige weitere Anstrengung gefallen lassen müssen, sich aber zugleich umso mehr dazu beglückwünschen, 1846 nicht die Hände in den Schoß gelegt zu haben. In beiden Fällen wird man aber die weiteren Maßnahmen im Bewusstsein ihrer Sicherheit oder ihrer Notwendigkeit treffen können.

Dass mit Annahme dieser Regelung der Anfang des §. 10 des Regulativs von selbst entfällt, der in seiner unklaren Fassung als eine Hauptquelle des jetzigen Übels betrachtet werden muss, braucht nicht eigens erwähnt zu werden.

Bei der Anwendung der in diesem Plan enthaltenen Regelung darf die Witwenkasse keinen
\SourceLeaf{166}
Unterschied zwischen den jetzigen Witwen und den künftig hinzukommenden machen; denn nur unter der Voraussetzung einer völlig gleichen Behandlung ist es möglich geworden, das Defizit auf den mäßigen Rest von 2963 Rthl. zu senken. Außerdem würden sich die neuen Witwen benachteiligt fühlen, wenn den früheren eine besondere Bevorzugung \textit{aus der Kasse} gewährt würde. Denn die Möglichkeit, bei so milden Änderungen der Pensionen zu bleiben, beruht allein auf der Beitragserhöhung, an der die Ehemänner der neuen Witwen teilgenommen haben werden, die der früheren aber nicht. Da es andererseits jedoch wünschenswert erscheint, jeder Unzufriedenheit der gegenwärtigen Witwen vorzubeugen, auch wenn sie nicht im Recht begründet wäre, schlägt die Kommission vor,
\begin{quote}
dass die Universität an das h. Kuratorium die Bitte richtet, den gegenwärtigen Witwen, soweit sie dann noch am Leben sind, für die Zeit, in der die Pension infolge der obigen Regelung unter 250 Rthl. sinkt, den an 250 Rthl. fehlenden Betrag aus der Universitätskasse ergänzen zu lassen.
\end{quote}
\noindent\hspace*{\parindent}Eine solche Bitte ließe sich mit naheliegenden Gründen unterstützen. Auch ist leicht zu erkennen, dass eine solche Beihilfe der Universitätskasse keine große Last auferlegen könnte. Es wäre sogar nicht unmöglich, dass der Fall einer Beihilfe \textit{niemals} einträte. Sollte sie aber schon bald erforderlich werden, weil eine neue Witwe hinzukommt, bevor eine der bisherigen ausgeschieden ist, so wird der Zuschuss für jede einzelne Witwe voraussichtlich doch für geraume Zeit nur wenige Taler betragen können. In späterer Zukunft, zum Beispiel nach 30 Jahren, könnte, wenn wir beispielhaft einen der äußersten Fälle annehmen, dass nämlich bis dahin überhaupt keine Pensionserhöhung zulässig geworden und die Zahl der Witwen auf 30 gestiegen wäre, der Zuschuss für eine Witwe 20 Rthl. betragen. Dann werden aber wahrscheinlich nur noch wenige der jetzigen Witwen am Leben sein.

Die Kommission glaubt, hier noch ein paar andere Einrichtungen erwähnen zu müssen, die als Mittel zur Verbesserung der Bilanz zur Sprache gekommen sind, ohne jedoch einen Antrag darauf richten zu wollen.

Zunächst ist die Frage aufgekommen, ob eine gewisse Verbesserung der Bilanz dadurch erreicht werden könnte, dass den künftig eintretenden außerordentlichen Professoren nur der Anspruch auf eine geringere Witwenpension eingeräumt würde. Sie würden damit eine zweite Klasse von Beteiligten bilden, aus der sie bei ihrer Beförderung auf eine ordentliche Professur von selbst in die erste aufrücken. Die Beiträge der Mitglieder zweiter Klasse sollten dagegen bei ihrer bisherigen Höhe von 2 Louis d'or bleiben. Den jetzigen außerordentlichen Professoren sollte die Wahl gelassen werden, ob sie in der ersten Klasse bleiben oder in die zweite übertreten wollten, jedoch so, dass im letzteren Fall ein Rücktritt in die erste Klasse nicht zulässig wäre, solange sie außerordentliche Professoren blieben.

Zur richtigen Würdigung dieses Gedankens ist zunächst sorgfältig zu bedenken, dass in dem oben S. [160] aufgestellten Hauptplan die Verringerung des Defizits von 17520 Rthl. auf 2963 Rthl. wesentlich davon abhängt, dass die Beiträge \textit{aller} Mitglieder, der jetzigen wie der künftigen, von 2 auf 3 Louis d'or erhöht werden. Dieser Hauptplan kann somit überhaupt nicht bestehen, wenn diese Voraussetzung eine wesentliche Einbuße erleidet, ohne dass dafür anderswo ein vollständiger und sicherer Ersatz eintritt. Damit ist eine Art, die Pension der zweiten Klasse niedriger festzusetzen, die sich sonst durch ihre Einfachheit empfehlen würde, von selbst als unzulässig ausgeschlossen: nämlich den Witwen zweiter Klasse nur den festen Teil der Pension, also 200 Rthl., zu gewähren oder sie völlig der Professorenwitwe G. gleichzustellen. Tatsächlich ist nämlich leicht zu erkennen, dass dadurch sehr wahrscheinlich \textit{die Kasse} überhaupt keinen Ersatz für jene Einbuße an Beiträgen erhalten würde. Nur die Witwen erster Klasse würden etwas besser gestellt, da die Zusatzdividende unter eine geringere Anzahl von Beteiligten verteilt würde.

